\subsection{An independent interval certificate}\label{push:section:E.4}
\label{sec:e4}

We write $\eta_D$ for the Dedekind eta function and set
\[
 h=z-(1-\tau)/2,\qquad
 \mathcal N^2(z,\tau)=
 \left|\frac{\vartheta_3(h,\tau)}{\eta_D(\tau)}\right|^2
 \exp\left(-\frac{2\pi(\Im h)^2}{\Im\tau}\right).
\]
The half-period theta identity identifies this with the norm used in
\cref{fam:src:C.1.1,fam:src:C.2.3}. The coefficient of $\Im c$ in $D_c$ is one. Consequently the
identity from the period calculation above which we use is
\[
 D_{\rm part}=\frac6\pi
 \log\left(\frac{\mathcal N}{2\sqrt3}\right).
\]
It follows directly by taking the imaginary part of the theta formula
for $\beta^{\rm part}$; its Gaussian factor removes
$6(\Im\mu)^2/\Im\tau$. Thus no numerical normalization is inferred
from the new computation. The analytic formula and the first
certificate are recorded in the period calculation above, \cref{fam:src:C.1.1,fam:src:C.2.3}. The result below supplies independent
arithmetic evidence for that same conclusion, with the independent certificate below.

\begin{theorem}[independent certification of every box]\label{push:src:E.4.1}
For all 19701 rational rectangles in the supplied certificate, and for
each $\tau$ in the rectangle, both zeros of $\wp_\tau$ satisfy
$\mathcal N^2<2$. The new largest certified upper bound is less than
\[
 1.999877571767410525.
\]
The 19701 rectangles, 36 local rectangles and 174 excluded rectangles
have disjoint interiors and cover exactly
$[0,1/2]\times[17/20,2]$. Every local rectangle is contained in
$|\tau-i|<3/20000$, and every excluded rectangle satisfies
$|\tau|<1$ throughout.
\end{theorem}
\Uses{the hypotheses, notation and definitions stated here; \citet[Sections 2--3]{Johansson2016}; \citet{PythonFlint2026}}


\begin{proof}
The programs \path{code/e4_arb_certificate.py} and
\path{code/e4_cover_and_local.py} use respectively 192-bit and 256-bit
Arb arithmetic through python-flint 0.9.0. The arithmetic implementation
is independent of mpmath. We import neither mpmath nor any routine
from the previous certificate. The only numerical data taken from that
certificate are its rational rectangle endpoints. We recompute roots,
preconditioners, radii, contraction bounds, function values and tails.
Arb's midpoint-radius representation and its outward error bounds are
described by \citet[Sections 2--3]{Johansson2016}; the precise
constructor and endpoint semantics used here are documented in
\citet{PythonFlint2026}. The final comparison is made with an upper
endpoint, before conversion to a decimal string.

For a rectangle $Q$, let $z_0$ and $A\ne0$ be the exact dyadic complex
numbers selected by the midpoint Newton iteration. That iteration is
only a choice of constants. For $Z=z_0+[-r,r]+i[-r,r]$ we independently
enclose $\wp_\tau(z_0)$ and $\wp'_\tau(Z)$ for $\tau\in Q$.
We accept only if
\[
 \begin{split}
 &\left|-A\wp_\tau(z_0)
 +(1-A\wp'_\tau(Z))(Z-z_0)\right|_\infty<r,\\
 &|\Re(1-A\wp'_\tau(Z))|
       +|\Im(1-A\wp'_\tau(Z))|<1.
 \end{split}
\]
Here absolute values on enclosures mean their outward upper bounds.
The derivative enclosure and integration along a line segment show
that $z\mapsto z-A\wp_\tau(z)$ is a strict contraction of $Z$ into
itself, uniformly for $\tau\in Q$. Thus it has a genuine zero in $Z$.
The two zeros of $\wp$ are negatives modulo the period lattice, counted
with multiplicity. The theta transformation laws show that
$\mathcal N$ is invariant under these operations; hence one root
enclosure controls both zeros. No continuity or success of Newton
iteration is used in that conclusion.

We use the normally convergent cosine expansion
\[
 \wp_\tau(z)=\frac{\pi^2}{\sin^2\pi z}-\frac{\pi^2}{3}
 +8\pi^2\sum_{n\geq1}\frac{nq^n}{1-q^n}
                      (1-\cos 2\pi nz),\qquad q=e^{2\pi i\tau},
\]
and differentiate it in $z$ for $\wp'$. This identity follows by twice
differentiating the normally convergent Jacobi product and inserting
$E_2=1-24\sum_{n\geq1}nq^n/(1-q^n)$; this also verifies the
Fourier formula used in the period calculation above. The new program sums
32 terms directly in cosine form. Put $r_q\geq|q|$,
$H_z\geq|\Im z|$, and $R_q=r_q e^{2\pi H_z}<1$.
For $j=1,2$, the inequality $(N+k)^j\leq(N+1)^jk^j$ gives
\[
 \begin{aligned}
 U_1(a)&=\frac{(N+1)a^{N+1}}{(1-a)^2},\\
 U_2(a)&=\frac{(N+1)^2a^{N+1}(1+a)}{(1-a)^3},
 \qquad
 \sum_{n>N}n^j a^n\leq U_j(a).
 \end{aligned}
\]
We therefore add complex rectangles enclosing discs of radii
\[
 \frac{8\pi^2}{1-r_q}\bigl(U_1(r_q)+U_1(R_q)\bigr),
 \qquad \frac{16\pi^3}{1-r_q}U_2(R_q)
\]
to the two partial sums. These estimates follow from
$|\cos z|,|\sin z|\leq e^{|\Im z|}$ and
$|1-q^n|\geq1-r_q$. They differ from the moment formulas used in
the old implementation.

For $\vartheta_3$ we sum
$1+2\sum_{n=1}^{12}e^{\pi i n^2\tau}\cos(2\pi nh)$.
With $T_*\leq\Im\tau$ and $H_*\geq|\Im h|$, the remaining
absolute value is at most
\[
 \frac{2e^{-169\pi T_*+26\pi H_*}}
 {1-e^{-27\pi T_*+2\pi H_*}}.
\]
The ratio in this expression is tested to be less than one. For eta
we independently sum its logarithm,
\[
 \log\eta_D(\tau)=\frac{\pi i\tau}{12}
          +\sum_{n=1}^{24}\log(1-q^n)+\mathcal E,
 \qquad
 |\mathcal E|\leq
 \frac{r_q^{25}}{(1-r_q)(1-r_q^{25})}.
\]
The logarithms are holomorphic on $|q^n|<1$. The tail follows by
summing $|\log(1-z)|\leq|z|/(1-|z|)$.
We enclose $\mathcal N^2$ as
\[
 ((\Re\vartheta_3)^2+(\Im\vartheta_3)^2)
 \exp\left(-2\Re\log\eta_D-
                   \frac{2\pi(\Im h)^2}{\Im\tau}\right),
\]
and test its upper endpoint against the exact integer 2.

All 19701 tests succeeded. The full output records each new enclosure;
the largest is enclosed by
\[
 1.999877571767410524070536478244687423578792301873142913
       \ \pm 9.26\cdot10^{-56}.
\]
This is the enclosure of a certified upper bound, not an estimate of
the actual maximum. Its box has input index 17866, with zero-based
indexing. The file digest of the input is recorded in the output.

The cover check is exact rational arithmetic. At all 438 distinct
vertical edges we maintain the active closed vertical intervals. In
every open vertical slab they form an ordered partition of
$[17/20,2]$, each with multiplicity one. The last active set is empty;
all leaves lie in the root rectangle, and the sum of their exact areas
is $23/40$. The closed rectangles also cover their common edges.
The omitted and excluded conditions are checked separately using
rational endpoint comparisons. Finally,
$2(1/10000)^2<(3/20000)^2$ puts the entire omitted square inside the
analytic disc used below. This verifies both the inequalities and
the coverage, rather than just replaying the old list of successful
flags.
\end{proof}

\begin{theorem}[quantitative local maximum]\label{push:src:E.4.2}
If $|\tau-i|\leq3/20000$ and $\wp_\tau(z)=0$, then
$\mathcal N(z,\tau)\leq\sqrt2$, with equality precisely at
$\tau=i$. The proof uses the explicit constants below.
\end{theorem}
\Uses{the hypotheses, notation and definitions stated here; the explicit constructions and calculations in the proof}


\begin{proof}
Set $x=\tau-i$, $G(h,\tau)=\wp_\tau((1-\tau)/2+h)$ and
$H(h,\tau)=\log(\vartheta_3(h,\tau)/\eta_D(\tau))$.
Both functions are even in $h$. On $|h|\leq1/5$, $|x|\leq1/10$,
the point $(1-\tau)/2+h$ has distance at least $0.9/2-0.2$ in
imaginary part from any lattice point, so $G$ has no pole.
Put
\[
 A_k=\sum_{n\geq1}n^ke^{-0.9\pi n^2}\cosh(0.4\pi n),
 \quad a_0=2A_0,\ a_1=4\pi A_1,\ a_2=8\pi^2A_2,
 \quad r=e^{-1.8\pi}.
\]
Then $a_0<0.225$, so the theta series has no zero there and a
holomorphic logarithm exists. From
$G=-\partial_h^2\log\vartheta_3-\pi^2E_2/3$ we obtain
\[
 \begin{aligned}
 |G|&\leq\frac{a_2}{1-a_0}+
 \frac{a_1^2}{(1-a_0)^2}+
 \frac{\pi^2}{3}\left(1+\frac{24r}{(1-r)^3}\right)<18.373<20,\\
 |H|&\leq-\log(1-a_0)+\frac{1.1\pi}{12}
                  +\frac{r}{(1-r)^2}<0.547<0.6.
 \end{aligned}
\]
The new certified upper bounds are respectively
$18.372504032305989$ and $0.546229724979279$.
For the sums $A_k$ we retain ten terms. If $n=11+j$, then
$n^2\geq121+23j$ and $n^k\leq11^k(1+j)^k$; replacing cosh by
the positive exponential bounds the tail by a geometric series with
ratio $e^{-20.3\pi}$ and its first two moment sums. These formulas,
including all outward rounding, are in the accompanying program.

At $(0,i)$ the coefficients are
\[
 G=0,\qquad G_{h^2}=-g_2(i)/4,\qquad G_x=id,
 \qquad H=\log\sqrt2,\quad H_x=0,\quad H_{h^2}=-\pi/2,
\]
where $d>0$. The first two identities follow from the half-period
value of $\wp$ and $\wp''=6\wp^2-g_2/2$. If
$e_2(\tau)=\pi^2(\vartheta_2^4-\vartheta_4^4)/3$, differentiating
its Gaussian series gives $e_2'(i)=id$. We compute
\[
 \begin{split}
 10.03&<d<10.04,\\
 47.26&<g_2(i)/4<47.30.
 \end{split}
\]
More precisely the certified enclosures begin
$d=10.03061933746602249879\ldots$ and
$g_2(i)/4=47.26818003230846307326\ldots$.
The computation of $d$ uses 16 terms in the two Gaussian series and
their derivatives, followed by geometric Gaussian tail bounds.
For $g_2/4$ it uses
\[
 \frac{g_2(i)}4=\frac{\pi^4}{3}
 \left(1+240\sum_{n\geq1}\frac{n^3e^{-2\pi n}}{1-e^{-2\pi n}}\right)
\]
with 24 terms and the bound
$\sum_{n>24}n^3q^n/(1-q^n)
\leq25^3q^{25}(1+4q+q^2)/(1-q)^5$.
The modular transformation of $\vartheta_3/\eta_D$ gives $H_x=0$;
that of eta gives $E_2(i)=3/\pi$. Inserting this in the displayed
formula for $G$ gives $H_{h^2}=-\pi/2$.

Take $\epsilon=3/20000$, $r_h=3/100$,
$s=\epsilon/(1/10)=3/2000$ and $r_2=(r_h/(1/5))^2=9/400$.
Cauchy's coefficient estimates give on $|h|=r_h$
\[
 |G(h,i)|\geq
 \left(47.26-\frac{500r_2}{1-r_2}\right)r_h^2
       >0.03217594373401534,
\]
while
\[
 |G(h,i+x)-G(h,i)|\leq
 \frac{20s}{(1-s)(1-r_2)}<0.03073664204747019.
\]
Rouch\'e's theorem therefore encloses both zeros in $|h|<r_h$.
The comparison with the total number of zeros of an elliptic function
of order two identifies these with all roots modulo the lattice.

Write $G=e_2(\tau)+h^2Q(h,\tau)$. The same Cauchy bounds imply
\[
 \left|Q+\frac{g_2(i)}4\right|
 \leq\frac{500(r_2+s/(1-s))}{1-r_2}<12.277368,
 \quad
 |e_2(\tau)-idx|\leq\frac{2000|x|^2}{1-s}.
\]
At a root, first bounding the left side of the last inequality below
and then bounding it above yields respectively
\[
 |x|<6.123343|h|^2<7|h|^2,
 \qquad |h|^2<0.000044338219<\frac1{20000}.
\]
For example the two rational majorants actually evaluated are
\[
 \frac{47.30+Q_*}{10.03-2000\epsilon/(1-s)},
 \qquad
 \frac{10.04\epsilon+2000\epsilon^2/(1-s)}{47.26-Q_*},
\]
where $Q_*=500(r_2+s/(1-s))/(1-r_2)$; both denominators are
strictly positive.

In the even power series of $H$, put
$\alpha=|x|/(1/10)\leq70|h|^2$ and
$\beta=|h|^2/(1/5)^2=25|h|^2$. After removal of its displayed
constant and linear coefficients, the remaining bivariate series is
bounded by
$0.6(\alpha^2+\alpha\beta+\beta^2)/((1-\alpha)(1-\beta))$.
The change from $1/\Im\tau$ to one in the Gaussian contributes
at most $7\pi|h|^4/(1-\epsilon)$. Thus a bound for the total
coefficient of $|h|^4$ is
\[
 \frac{0.6(70^2+25\cdot70+25^2)}
 {(1-s)(1-25/20000)}+\frac{7\pi}{1-\epsilon}
 <4399.024<5000.
\]
Since $\Re(-\pi h^2/2)-\pi(\Im h)^2=-\pi|h|^2/2$, we have
\[
 \log\mathcal N-\log\sqrt2
 \leq-\frac\pi2|h|^2+5000|h|^4
 <-\left(\frac\pi2-\frac14\right)|h|^2
\]
whenever $h\ne0$. The coefficient is greater than $1.32079632679$.
If $h=0$, the lower bound for $e_2$ forces $x=0$. This proves the
claim with all constants controlled by a second arithmetic engine.
\end{proof}

\paragraph{Completion of the independent proof of Theorem~\ref{fam:st:C.2.3}.}

\begin{proof}
The theta transformation formulas preserve $\mathcal N$ under modular
transport, lattice translation, sign, and
$(z,\tau)\mapsto(\bar z,-\bar\tau)$; hence the half modular domain
suffices. Theorems \ref{push:src:E.4.1}--\ref{push:src:E.4.2} cover its portion $\Im\tau\leq2$.
For $T=\Im\tau\geq2$, choose $|\Im z|\leq T/2$ by lattice
translation. With $b=43/100$ and $r=e^{-4\pi}$, if
$|\Im z|\geq b$ the Fourier formula gives
\[
 \left|\frac{\wp_\tau(z)}{\pi^2}+\frac13\right|
 \leq\frac1{\sinh^2(\pi b)}+
 8\left\{\frac r{(1-r)^3}
 +\frac1{2(1-r)}\left(
 \frac{e^{-2\pi}}{(1-e^{-2\pi})^2}+
 \frac{e^{-6\pi}}{(1-e^{-6\pi})^2}\right)\right\}<\frac13.
\]
Thus a zero has $|\Im z|<b$. The function
$\cosh(\pi y)e^{-\pi y^2/T}$ is increasing for $0\leq y\leq b$:
concavity and the positive endpoint value of $\tanh(\pi y)-y$
give its nonnegative logarithmic derivative. The Jacobi product now
implies
\[
 \mathcal N\leq2\cosh(\pi b)
 e^{-\pi/3-\pi b^2/2}
 \exp\left(\frac{2\cosh(2\pi b)r}{1-r}
                +\frac{r^2}{1-r^2}\right)<1.082<\sqrt2.
\]
This bound decreases with $T\geq2$. Both strict inequalities are
independently checked in \path{code/e4_cover_and_local.py}.
At $\tau=i$, $\vartheta_3(0,i)/\eta_D(i)=\sqrt2$, so equality
does occur. The period identity at the beginning of the subsection
then gives the claimed value of $M$ and
$e^{-\pi M}=6^3=216$. The invariant $j(\tau)=1728t$ sends its
unique modular equality orbit to $t=1$.
\end{proof}



\clearpage\appendix
